\section{Weekly Summary}

\begin{tcolorbox}[colback=yellow!10,colframe=black!50,title=AI-Generated Content]
\noindent The summary below was written by Claude (Anthropic) from that week's regression outputs and series descriptions. It has not been reviewed or fact-checked by a human, and should be read as an automated, informal recap -- not analysis.
\end{tcolorbox}

Welcome back to the only recurring column on the internet where we throw two completely unrelated economic time series into a blender and act deeply moved by whatever comes out. This week the machine served up a spicy inflation index against a discontinued interest rate benchmark, a defunct Swiss central bank intervention log against manufacturers' order ratios, and, for the grand finale, Asian air freight import prices against how many people were hiring for burger-flipping jobs on Indeed. Truly the intellectual heirs of Newton and Keynes are shaking in their boots.

Monday, August 10, gave us our headline act: the Chained CPI for Services versus the 1-Week AMERIBOR term structure, and oh boy, the raw regression came out swinging with an R-squared of 0.93. Ninety-three percent! Somebody call Stockholm! Except, as any regular reader of this dunk-tank of a newsletter already suspects, both series just spent their overlapping lives drifting upward like two balloons let go at the same birthday party. Once we detrend and strip that shared drift out, the R-squared collapses to a much more humble 0.23. Now, to be fair, that detrended relationship is still technically "significant" at the usual thresholds, which is more than most of our contestants can say -- so I'll grudgingly admit this one didn't completely evaporate. It went from "revolutionary discovery" to "eh, maybe there's a faint something," which in this project counts as a rousing success. Do not, under any circumstances, build a portfolio on it.

Wednesday, August 12, was the day the math simply gave up. Pairing decades-old Swiss National Bank USD-against-DEM purchases with a transportation-equipment order ratio produced an R-squared of roughly 0.0000000000000003, a slope of exactly zero, and a parade of p-values that just said "nan" like a confused toddler. The design matrix was reportedly singular, which is a polite statistical way of saying the two series had absolutely nothing to say to each other and the software knew it. A moment of silence for the numbers that died so this row of the CSV could exist. Then Thursday's air-freight-versus-fast-food-job-postings matchup limped in with an R-squared of 0.02 raw and a p-value hovering around a deeply unconvincing 0.19, meaning it wasn't significant in the raw version and stayed just as unimpressive after detrending. Consistency! Just... consistently nothing.

As always, the mandatory reminder for anyone who wandered in expecting rigor: these pairs are chosen at random, nobody is claiming Swiss currency interventions cause anything, and spurious correlation isn't a bug here, it's the entire genre. The one takeaway worth keeping is the meta-lesson the CPI/AMERIBOR pair illustrates so nicely -- a jaw-dropping 0.93 in levels means almost nothing on its own, and the only reason we even raise an eyebrow at that pair is that a sliver of it survived detrending. Everything else this week was, statistically speaking, two strangers waving at each other from different trains. See you next week for more.
